\documentclass{situationist-preprint}

\title{Ordinary Objects Vindicated}
\msid{MET-06-c}
\msfamily{MET-06}
\msdiscipline{Metaphysics}
\mstype{Consequence}
\mskeywords{ordinary objects; common fate; debunking arguments; object perception; special composition question}

\begin{document}
\maketitle

\begin{abstract}
MET-06-a answers the special composition question with the common-fate criterion and derives that tables and organisms exist while arbitrary scattered sums do not. We draw the consequence for ordinary objects in general. First, the criterion sorts the objects of common sense into those that exist as composites and those, such as heaps and flocks, that are pluralities about which ordinary sentences are true without a further object. Second, it answers the debunking argument that our beliefs about which objects exist are formed by processes insensitive to the facts about composition. Perceptual object segregation is driven by common motion, which is the evidence of common fate, so ordinary object beliefs track the very condition that, by the criterion, makes them true.
\end{abstract}

\section{Which ordinary objects exist}

By Theorem~4.1 of MET-06-a, any composition rule satisfying Unity, Parity and Ordinary Objects coincides with the common-fate rule: some simples compose something exactly when, under every admissible perturbation of their environment, they have the same fate. Corollary~4.3 applies this to tables, organisms and scattered pluralities. Applied to the inventory of common sense, the rule yields the following division.

\paragraph{Composites.} Rigid artefacts (tables, hammers, cars), organisms, stones, molecules, planets, and the cohesive undetached parts of each, such as the legs of a table. In each case the members are bound so that a push on one is a push on all.

\paragraph{Pluralities.} Heaps of sand, flocks, crowds, forests and piles of stored planks. Their members are correlated in position or motion, but an admissible perturbation of one member does not carry the others with it (MET-06-a, Section~6). There is no further object, but ordinary sentences about them are true: \enquote{the heap is in the corner} is true because the grains are in the corner, and \enquote{the flock turned south} because the birds turned south together. These are plural predications, and the criterion leaves them intact.

\paragraph{Borderline cases.} Cairns, liquids in containers, and loosely joined assemblies. On the epistemic view, which MET-06-a adopts from MET-18, each such case falls on one side or the other of a sharp but unknown boundary (MET-06-a, Corollary~7.2).

This division coincides with common sense in its verdicts about existence where common sense is confident, and departs from it only in ways common sense itself marks. Nobody thinks a flock is a body in the way a bird is, and the habit of speaking of \enquote{a heap} is matched by the habit of speaking of \enquote{some sand}. The vindication of ordinary objects is therefore not wholesale: it is the vindication of the ordinary distinction between things and collections.

\section{The debunking argument}

\citet{Korman2015} identifies the strongest argument against ordinary objects as a debunking argument. Our beliefs about which objects there are, the argument goes, are produced by perceptual and conceptual processes shaped by evolutionary and cultural pressures that have nothing to do with the facts about composition. Had those pressures been different, we would have believed in trout-turkeys and not in trout. Since the beliefs would have been held whether or not they were true, they are not knowledge. Combined with the eliminativist arguments of \citet{Merricks2001}, the conclusion would be that we have no reason to believe in ordinary inanimate objects at all.

\begin{proposition}\label{p:tracking}
If composition is common fate, the processes that produce our beliefs about which material objects exist are sensitive to the facts that make those beliefs true.
\end{proposition}

\begin{proof}
The perceptual systems that segregate the visual scene into objects group regions that move together and separate regions that move independently. Common motion is the grouping principle that \citet{Wertheimer1923} called the law of common fate, and in infants it is among the earliest and most robust cues to object unity: surfaces that move rigidly together are perceived as one object, and surfaces that move apart as two \citep{Spelke1990}. Moving together under perturbation is the observable manifestation of sharing a common fate in the sense of MET-06-a (Definition~3.4). So the process that produces the belief that some region is one object responds to the condition that, by the common-fate criterion, makes the belief true: had the parts not been bound, they would have come apart under the ordinary perturbations of the environment, and the process would have represented them as many.
\end{proof}

Proposition~\ref{p:tracking} does not say that perception is infallible. Two stones resting against each other look like one until one is pushed, and a camouflaged animal is not seen as distinct from its background until it moves. But the errors are errors about whether a perturbation has revealed the binding, and they are corrected by the same process when it does. The debunker's counterfactual, that we would have believed in trout-turkeys under different pressures, is true only of creatures whose object perception did not group by common motion; and such creatures would have been tracking something other than the conditions under which material things hang together.

\begin{objection}[Coincidence of criterion and perception]
The fit between perception and the criterion is suspicious. The criterion was chosen to match ordinary judgments, and the judgments are the product of perception; so the fit is guaranteed, not discovered.
\end{objection}

\begin{reply}
The criterion was not chosen to match judgments but derived from principles, Unity and Parity, together with a single Moorean premise about tables (MET-06-a, Section~5). Unity is motivated by what a composite is for, the unity of its career under perturbation, and not by perceptual grouping. That a principle so motivated coincides with the grouping principle of the visual system is evidence that both respond to the same feature of the world, and it is the coincidence a non-sceptical account of object perception would predict.
\end{reply}

\section{Conclusion}

The common-fate criterion vindicates the objects of common sense where common sense treats something as a body, and it explains the remainder as plural reference to things that are not bound. It also explains why our beliefs about ordinary objects deserve trust: the cue by which we perceive things as one is the observable sign of the condition under which they are one.

\begin{thebibliography}{9}\small
\bibitem[Korman(2015)]{Korman2015} Korman, D.~Z. (2015). \emph{Objects: Nothing Out of the Ordinary}. Oxford University Press.
\bibitem[Merricks(2001)]{Merricks2001} Merricks, T. (2001). \emph{Objects and Persons}. Oxford University Press.
\bibitem[Spelke(1990)]{Spelke1990} Spelke, E.~S. (1990). Principles of object perception. \emph{Cognitive Science}, 14(1), 29--56.
\bibitem[Wertheimer(1923)]{Wertheimer1923} Wertheimer, M. (1923). Untersuchungen zur Lehre von der Gestalt, II. \emph{Psychologische Forschung}, 4, 301--350.
\end{thebibliography}

\end{document}
